\chapter{Controller Design and Testing}
\label{ch:controller-design}

\epigraph{A controller is only as good as its worst-case performance.}{}

\section{PD Control with Gravity Compensation}

\begin{lstlisting}[language=Python, caption={PD controller with gravity compensation.}]
import numpy as np

class PDGravityCompController:
    """PD controller with gravity compensation.

    From Volume I, Chapter 3: the simplest nonlinear controller
    that guarantees asymptotic stability.
    """

    def __init__(self, Kp: np.ndarray, Kd: np.ndarray,
                 gravity_func=None):
        self.Kp = Kp
        self.Kd = Kd
        self.gravity_func = gravity_func

    def compute(self, q: np.ndarray, qd: np.ndarray,
                q_des: np.ndarray, qd_des: np.ndarray) -> np.ndarray:
        """Compute control torques.

        Parameters
        ----------
        q, qd : current state
        q_des, qd_des : desired state

        Returns
        -------
        tau : np.ndarray, control torques
        """
        tau = self.Kp @ (q_des - q) + self.Kd @ (qd_des - qd)
        if self.gravity_func is not None:
            tau += self.gravity_func(q)
        return tau
\end{lstlisting}

\section{Trajectory Tracking with Contraction}

Implementing the contraction-based controller from Volume~I, Chapter~4.

\section{Funnel Control}

Implementing trajectory tubes (Volume~II, Chapter~3) for robust tracking.

\section*{Exercises}
\begin{enumerate}
  \item Implement PD control for the 2-DOF arm and track a circular trajectory.
  \item Add gravity compensation and compare tracking error.
  \item Implement a contraction-based controller and verify convergence.
\end{enumerate}
